% Part: propositional-logic
% Chapter: syntax-and-semantics
% Section: preliminaries

\documentclass[../../../include/open-logic-section]{subfiles}

\begin{document}

\olfileid[tr]{pl}{syn}{pre}

\olsection{Ön Bilgiler}

\begin{thm}[\emph{!!{formula}s üzerinde tümevarım ilkesi}]
  \ollabel{thm:induction}  
  Bir $P$ özelliği bütün atomik !!{formula}s için geçerliyse ve
  \begin{enumerate}
    \tagitem{prvNot}{$!A$ için geçerli olduğunda $\lnot !A$ için de
      geçerliyse;}{}
    \tagitem{prvAnd}{$!A$ ve~$!B$ için geçerli olduğunda $(!A \land !B)$
      için de geçerliyse;}{}
    \tagitem{prvOr}{$!A$ ve~$!B$ için geçerli olduğunda $(!A \lor !B)$
      için de geçerliyse;}{}
    \tagitem{prvIf}{$!A$ ve~$!B$ için geçerli olduğunda $(!A \lif !B)$
      için de geçerliyse;}{}
    \tagitem{prvIff}{$!A$ ve~$!B$ için geçerli olduğunda $(!A \liff !B)$
      için de geçerliyse;}{}
  \end{enumerate}
  o hâlde $P$ bütün !!{formula}s için geçerlidir.
\end{thm}

\begin{proof}
  $S$, $P$ özelliğine sahip bütün !!{formula}s topluluğu olsun. Açıkça
  $S \subseteq \Frm[L_0]$. $S$, \olref[fml]{defn:formulas}'daki bütün
  koşulları sağlar: bütün atomik !!{formula}s kümesinin tamamını içerir ve
  !!{operator}s
  altında kapalıdır. $\Frm[L_0]$ bu türden en küçük sınıftır; dolayısıyla
  $\Frm[L_0] \subseteq S$. Öyleyse $\Frm[L_0] = S$ ve her formül $P$
  özelliğine sahiptir.
\end{proof}

\begin{prop}\ollabel{prop:balanced}
   $\Frm[L_0]$'daki her !!{formula}, sol parantezlerinin sayısı sağ
   parantezlerinin sayısına eşit olması bakımından \emph{dengelidir}.
\end{prop}

\begin{prob}
\olref[pl][syn][pre]{prop:balanced}'i kanıtlayın.
\end{prob}

\begin{prop} \ollabel{prop:noinit}
!!a{formula} verilsin. Bu formülün hiçbir öz başlangıç parçası başka bir
!!a{formula} değildir.
\end{prop}

\begin{prob}
  \olref[pl][syn][pre]{prop:noinit}'i kanıtlayın.
\end{prob}

\begin{prop}[Tek Türlü Okunabilirlik]
$\Frm[L_0]$'daki her !!{formula}~$!A$, aşağıdakilerden biri olarak tam bir
ayrıştırmaya sahiptir:
\begin{enumerate}
\tagitem{prvFalse}{$\lfalse$.}{}

\tagitem{prvTrue}{$\ltrue$.}{}

\item Bir $\Obj p_n \in \PVar$ için $\Obj p_n$.
  
\tagitem{prvNot}{Bir !!{formula}~$!B$ için $\lnot !B$.}{}

\tagitem{prvAnd}{Bazı !!{formula}s $!B$ ve~$!C$ için $(!B \land !C)$.}{}

\tagitem{prvOr}{Bazı !!{formula}s $!B$ ve~$!C$ için $(!B \lor !C)$.}{}

\tagitem{prvIf}{Bazı !!{formula}s $!B$ ve~$!C$ için $(!B \lif !C)$.}{}

\tagitem{prvIff}{Bazı !!{formula}s $!B$ ve~$!C$ için $(!B \liff !C)$.}{}
\end{enumerate}
Üstelik bu ayrıştırma \emph{tektir}.
\end{prop}

\begin{proof}
$!A$ üzerinde tümevarımla. Örneğin $!A$'nın $(!B \lif !C)$ ve
$(!B' \lif !C')$ olarak iki farklı okuması bulunduğunu varsayalım. O zaman
$!B$ ile $!B'$ aynı olmalıdır (yoksa biri diğerinin öz başlangıç parçası
olurdu); dolayısıyla $!A$'nın iki okuması farklıysa bunun nedeni, $!C$ ile
$!C'$'nin aynı simge dizisinin farklı okumaları olmasıdır; bu ise tümevarım
varsayımı uyarınca olanaksızdır.
\end{proof}


\begin{defn}[Düzgün Yerine Koyma]
$!A$ ve $!B$ !!{formula}s, $\Obj p_i$ de bir !!{propositional
variable}
ise $\Subst{!A}{!B}{\Obj p_i}$, $!A$'daki her $\Obj p_i$ geçişinin bir $!B$
geçişiyle değiştirilmesinin sonucunu gösterir; benzer biçimde $\Obj p_1$,
\dots,~$\Obj p_n$ yerine eşzamanlı olarak !!{formula}s $!B_1$,
\dots,~$!B_n$ konması
$\SSubst{!A}{\subst{!B_1}{\Obj p_1},\dots,\subst{!B_n}{\Obj p_n}}$
ile gösterilir.
\end{defn}

\begin{prob} Aşağıdaki listede beş tane !!{formula} vardır. Her biri için bu
 !!{formula}, \( \Subst{!A}{!B}{\Obj p_i} \) biçiminde ifade edilebilir mi,
 belirleyin; burada \( !A \), (i)~\( \Obj p_0 \); (ii)~\( ( \lnot \Obj p_0
 \land \Obj p_1) \); ve (iii)~\( ( ( \lnot \Obj p_0 \lif \Obj p_1 ) \land
 \Obj p_2 ) \) olsun. Her durumda ilgili yerine koymayı belirtin.
  \begin{enumerate}
    \item \( \Obj p_1 \) 
    \item \( ( \lnot \Obj p_0 \land \Obj p_0 ) \)
    \item \( ( ( \Obj p_0 \lor \Obj p_1 ) \land \Obj p_2 )  \)
    \item \( \lnot ( ( \Obj p_0 \lif \Obj p_1 ) \land \Obj p_2 ) \)
    \item \( (( \lnot ( \Obj p_0 \lif \Obj p_1 ) \lif ( \Obj p_0 \lor \Obj p_1 )) \land \lnot ( \Obj p_0 \land \Obj p_1 )) \)
  \end{enumerate}
\end{prob}

\begin{prob}
$\Subst{!A}{!B}{p}$ için tümevarımla matematiksel bakımdan kesin bir tanım
verin.
\end{prob}

\end{document}
